\input{algebra.tex}
\title[Algebra1.]{Polinomok}
\subtitle{Oszthatóság és maradékos osztás polinomgyűrűben}
\date[2. hét]{2026.~szeptember 20-án}
\begin{document}
\openingframes{Tartalom}
\section{Polinomgyűrűk}
\frame{
\begin{definition}[polinom]\index{polinom}
	Legyen $\mathbb{F}$ egy test.
	E test feletti polinomokon az összes
	\[
		p\left( t \right)=
		\alpha_0+\alpha_1t+\alpha_2t^2+\dots+\alpha_nt^n
	\]
	alakú formális algebrai kifejezést értjük.
	Itt $n$ tetszőleges nem negatív egész
	és $\alpha_0,\dots,\alpha_n$ tetszőleges, az $\mathbb{F}$ testbeli elemek.
	Az $\mathbb{F}$ test feletti összes polinomok halmazát $\mathbb{F}\left[ t \right]$ módon jelöljük.
\end{definition}
A fenti definícióban az \emph{algebrai kifejezés} szó arra utal,  hogy az
\begin{math}
	\alpha_0+\alpha_1t+\alpha_2t^2+\dots+\alpha_nt^n
\end{math}
műveletek minden $t\in\mathbb{F}$ mellett értelmesek,
és eredményük egy újabb $\mathbb{F}$ testbeli elem.
Ha $t\in\mathbb{F}$ konkrétan meg van adva,
akkor a behelyettesítés után kapott elemet mondjuk a $p$ polinom helyettesítési értékének.
}
\frame{
A \emph{formális algebrai kifejezés}\index{formális algebrai kifejezés} arra utal,
hogy egy polinomot az együtthatói határozzák meg,
azaz két polinom akkor és csak akkor azonos,
ha az megfelelő együtthatói azonosak.
Ez szemben áll avval, hogy ha a polinomokra mint függvényekre tekintenénk,
akkor a helyettesítési értékek egyenlősége jelentené a két polinom azonos voltát.
A formális szó tehát azt jelenti, hogy nem mint függvényre gondolunk,
hanem egyszerűen az adott $\alpha_0,\dots,\alpha_n$ rögzített elemek -- ezeket mondjuk együtthatóknak --,
által előírt műveletekre.
Az az előírás ugyanis, hogy tetszőleges $t\in\mathbb{F}$ mellett hajtsuk végre az
\[
	\alpha_0+\alpha_1t+\alpha_2t^2+\dots+\alpha_nt^n
\]
műveletsort.
A műveletsorról és nem annak eredményéről van szó.


Két műveletsor akkor azonos, ha ugyanazok a műveletsort meghatározó
$\left( \alpha_0,\alpha_1,\dots,\alpha_n \right)$
együtthatók.
}
\frame{
A jelölések megértése is fontos.
$p\left( t \right)\in\mathbb{F}\left[ t \right]$ semmi mást nem jelent,
minthogy $p\left( t \right)$ egy polinom.
Persze a polinom nem keverendő össze a helyettesítési értékével,
hiszen az egyik egy algebrai kifejezés-együttes, a másik egy az adott testbeli elem.
Szokásos viszont, hogy ha nincs konkrét $t$ a szövegkörnyezetben, akkor is $p\left( t \right)$ jelöli a polinomot.
Néha egyszerűbben csak $p$-vel jelöljük, főleg akkor ha nincs szó behelyettesítésről,
emiatt érdektelen a változó jele.
Ritkábban, de előfordul, hogy egy konkrét értékre, mondjuk $s\in\mathbb{F}$-re kell kiértékelnünk a polinomot ilyenkor $p\left( s \right)$ jelöli azt a testbeli elemet,
amelyet $t$ helyett $s$-et téve az előírt műveletek kiértékelése után kapunk.
A szövegkörnyezetben mindig világosnak kell lennie, hogy $p\left( t \right)$ a polinomot jelenti,
vagy egy konkrét $t$-re kiértékelt testbeli elemet.
}
\frame{\frametitle{Polinom fokszáma}
\begin{definition}\index{polinom foka}
	Legyen $p\left( t \right)\in\mathbb{F}\left[ t \right]$ egy polinom.
	Azt mondjuk, hogy az $n$ nem negatív egész szám e \emph{polinom fokszáma},
	ha $n$ a legnagyobb nemzérus együttható indexe.
	A legnagyobb nemzérus együtthatót \emph{főegyütthatónak}\index{főegyüttható} nevezzük.
	Azt mondjuk, hogy egy nemzérus polinom \emph{normált}\index{normált polinom}, ha $1$ a főegyütthatója.

	A $p\left( t \right)=0$ konstans zérus polinom foka megállapodás szerint legyen $-\infty$.
	A $p$ polinom fokszámát $\deg p$ módon jelöljük.
\end{definition}
Látni fogjuk, hogy a konstans zérus polinomra $\deg p=-\infty$ csak egy kényelmes jelölés.
Időnként a polinom fokszámával műveleteket is végzünk.
Megegyezés szerint ilyenkor $-\infty+a=-\infty$ minden $a$ nem negatív egész számra,
és $\left( -\infty \right)+\left( -\infty \right)=-\infty.$
A $-\infty$ szimbólumot minden egész számnál határozottan kisebbnek gondoljuk.
}

\frame{\frametitle{Polinomok műveletei}
Két polinom összegét és szorzatát a szokásos módon definiáljuk:
\begin{definition}\label{def:polmuveletek}
    Legyen $p,q\in\mathbb{F}[t]$, két polinom.
	\(
		p\left( t \right)
		=
		\sum_{j=0}^n\alpha_jt^j
		\text{ és }
		q\left( t \right)
		=
		\sum_{j=0}^m\beta_jt^j,
        \text{ itt }
		\alpha_j,\beta_j\in\mathbb{F},
		0\leq n,m\in\mathbb{Z}.
	\)

	Ekkor a $p$ és $q$ összegének és szorzatának definíciója:
	\[
		\left( p+q \right)\left( t \right)
		=
		\sum_{j=0}^{\max{\left\{ m,n \right\}}}\left( \alpha_j+\beta_j \right)t^j;
        \qquad
		\left( pq \right)\left( t \right)
		=
		\sum_{j=0}^{n+m}c_jt^j,
		\text{ ahol }
		c_j
		=
		\sum_{k=0}^j\alpha_k\beta_{j-k}.
		\qedhere
    \]
\end{definition}
Mind az összeadás, mind a szorzás definíciójában úgy kell a formulát érteni,
hogy amikor a hivatkozott együtthatók nem léteznek,
akkor értékük legyen zérus.
Például a szorzat legmagasabb indexű együtthatójára
$c_{n+m}=\alpha_n\beta_m$, hiszen az $\alpha_k,\beta_{n+m-k}$ számok
csak egyetlen $k$ mellett értelmezettek közösen, mikor $k=n$.
}
\frame{
	Ha már követjük a konvenciót,
	amely szerint a nem definiált együtthatókat zérusnak tekintjük,
	akkor persze
	\[
	\left( p+q \right)\left( t \right)
	=
	\sum_{j=0}^\infty\left( \alpha_j+\beta_j \right)t^j
    \text{ és }
	\left( pq \right)\left( t \right)
	=
	\sum_{j=0}^\infty\left( \sum_{k=0}^\infty\alpha_k\beta_{j-k} \right)t^j
    \]
	is írható.

	Itt a fenti összegek nem végtelen összegek, hiszen a két polinom együtthatói
	csak véges sok esetben különböznek zérustól.
	Így ebben az esetben a formálisan végtelen összeg definíciója
	egyszerűen a véges sok nem zérus elem összege.
}
\frame{
Fontos azt értenünk, hogy a $p$ és a $q$ polinom összegét és szorzatát úgy definiáltuk, hogy minden $s\in\mathbb{F}$ elemre a
\[
    \left( p+q \right)\left( s \right)
    =
    p\left( s \right)+q\left( s \right)
    \text{ valamint a }
    \left( p\cdot q \right)\left( s \right)
    =
    p\left( s \right)\cdot q\left( s \right)
\]
kiértékelések fennálljanak.
}
\frame{
\begin{proposition}
	Legyenek $p,q\in\mathbb{F}[t]$ polinomok az $\mathbb{F}$ test felett.
	Ekkor
	\begin{enumerate}
		\item $\deg \left( pq \right)=\deg p+\deg q$;
		\item $\deg \left( p+q \right)\leq\max\left\{ \deg p,\deg q \right\}$.\qedhere
	\end{enumerate}
\end{proposition}
%begin{proof}
%   Figyeljünk arra, hogy a konstans zérus polinom esetében is működik a tétel,
%   és vegyük észre, hogy az szorzat polinomra vonatkozó állítás azért igaz,
%   mert a test nullosztómentes.
%end{proof}
A következő állítás igazolása házi feladat.
Aprólékosan igazoljuk a test axiómák gyakorlásaként.
Vigyázat: a szorzás asszociativitása nem is olyan egyszerű.
\begin{proposition}
	Egy $\mathbb{F}$ test feletti $\mathbb{F}\left[ t \right]$ formális polinomok
	a fent bevezetett összeadás és szorzás műveletekkel,
	nullosztómentes,
	kommutatív, egységelemes gyűrűt alkotnak.
\end{proposition}
}
\section{Polinomok oszthatósága és a maradékos osztás}
\frame{\frametitle{Oszthatóság}
\begin{definition}\index{polinom osztója}
	Azt mondjuk, hogy a $p\in\mathbb{F}\left[ t \right]$ \emph{osztója} az $f\in\mathbb{F}\left[ t \right]$ nem zérus polinomnak,
	ha létezik $h\in\mathbb{F}\left[ t \right]$, hogy $f\left( t \right)=p\left( t \right)h\left( t \right)$.
	Ilyenkor $f$-et a $p$ egy \emph{többszörösének}\index{polinom többszöröse} is mondjuk.
	Jelölés: $p|f$.
\end{definition}
Világos, hogy egy $p$ polinom összes többszörösei -- tehát azok, amelyeknek $p$ osztója --
ideált alkotnak.
Ez a $p$ generálta legszűkebb ideál, azaz  a $J(p)=\left\{ fp:f\in\mathbb{F}\left[ t \right] \right\}$ főideál.
Ha $q\in J\left( p \right)$, akkor $J\left( q \right)\subseteq J\left( p \right)$, azaz ha $q$ egy többszöröse $p$-nek,
akkor $q$ minden többszöröse $p$-nek is többszöröse.
}

\metroset{block=fill}
\frame{
    \begin{block}{}
    Ha $p,q$ nem zérus polinomok,
    amelyekre $p|q$ és $q|p$ akkor a két polinom csak konstans szorzóban különbözik egymástól.
    \end{block}

Ha például a két polinom még normált is, akkor $p|q$ és $q|p$ csak $p=q$ esetben lehetséges.

A polinomgyűrű ideáljaira fokuszálva, azt gondoltuk éppen meg,
hogy \emph{$J\left( p \right)=J\left( q \right)$ normált $p,q$ polinomokra csak úgy teljesülhet, ha $p=q$},
azaz a polinomok gyűrűjében minden főideálnak csak egy generáló eleme van a normált polinomok körében.
}
\metroset{block=transparent}
\frame{
A következő állítás szerint a polinomok közt is működik a maradékos osztás,
ahogyan azt az egész számok közt megszoktuk.
\begin{proposition}[maradékos osztás]
	Legyenek $p,q\in\mathbb{F}\left[ t \right]$ polinomok, $q\neq 0$.
	Ekkor létezik egyetlen $h,r\in\mathbb{F}\left[ t \right]$ polinom, amelyre
	\[
		p
		=
		hq+r;
		\quad
		\deg r < \deg q.\qedhere
	\]
\end{proposition}
%begin{proof}
%   Először is azt vegyük észre, hogy $\deg p<\deg q$ esetben $r=p$,
%   $h=0$ szereposztással készen is vagyunk.
%
%   Tegyük fel tehát, hogy $n=\deg p\geq \deg q=m$, és lássuk be az állítást $n$ szerinti indukcióval.
%   Ha $n=0$, akkor $p\left( t \right)=\alpha_0$ és $q\left( t \right)=\beta_0\neq 0$.
%   Ekkor persze
%   \[
%   	\alpha_0=\frac{\alpha_0}{\beta_0}\beta_0+0,
%   \]
%   ami azt jelenti, hogy $h\left( t \right)=\frac{\alpha_0}{\beta_0}$ és $r\left( t \right)=0$ szereposztás
%   megfelelő.
%
%   Most tegyük fel,
%   hogy igaz az állítás $n+1$-nél kisebb fokú $p$ polinomokra ($n\geq 0$),
%   és lássuk be egy pontosan $n+1$-ed fokú polinomra.
%   Legyen tehát
%   \[
%   	p\left( t \right)=\alpha_{n+1}t^{n+1}+\dots+\alpha_0
%   	\quad\text{ és }\quad
%   	q\left( t \right)=\beta_{m}t^m+\dots+\beta_0,
%   \]
%   ahol $m\leq n+1$.
%   Tekintsük a következő polinomot:
%   \[
%   	\frac{\alpha_{n+1}}{\beta_m}t^{n+1-m}q\left( t \right).
%   \]
%   Világos, hogy ennek főegyütthatója éppen $\alpha_{n+1}$ és foka éppen $n+1=\deg p$.
%   Így a
%   \[
%   	p_1\left( t \right)
%   	=
%   	p\left( t \right)-
%   	\frac{\alpha_{n+1}}{\beta_m}t^{n+1-m}q\left( t \right).
%   \]
%   polinomra $\deg p_1<\deg p$.
%   Na most, ha $\deg p_1<\deg q$, akkor a bizonyítás első mondatában említett helyzetben vagyunk,
%   tehát nyilvánvaló szereposztással az állítás igaz $p_1$-re és $q$-ra.
%   Ha viszont $\deg p_1\geq \deg q$ még mindig igaz, akkor az indukciós feltétel szerint található
%   $h,r\in\mathbb{F}\left[ t \right]$ polinom, amelyre igaz az állítás.
%   Mindkét esetben találtunk tehát $h,r$ polinomokat, amelyre
%   \[
%   	p\left( t \right)-
%   	\frac{\alpha_{n+1}}{\beta_m}t^{n+1-m}q\left( t \right)
%   	=
%   	p_1\left( t \right)
%   	=
%   	h\left( t \right)q\left( t \right)+r\left( t \right);
%   	\quad
%   	\deg r < \deg q
%   \]
%   teljesül.
%   Ekkor persze
%   \[
%   	p\left( t \right)
%   	=
%   	\left( h\left( t \right)
%   	+
%   	\frac{\alpha_{n+1}}{\beta_m}t^{n+1-m}
%   	\right)
%   	q\left( t \right)
%   	+
%   	r\left( t \right);
%   	\quad
%   	\deg r < \deg q
%   \]
%   is fennáll. Ezt kellett belátni az állítás egzisztencia részéhez.
%
%   Az unicitás részhez tegyük fel, hogy valamely $h,h_1,r,r_1$ polinomokra
%   \[
%   	h\left( t \right)q\left( t \right)+r\left( t \right)
%   	=
%   	p\left( t \right)
%   	=
%   	h_1\left( t \right)q\left( t \right)+r_1\left( t \right)
%   \]
%   teljesül, ahol $\deg r<\deg q$ és $\deg r_1<\deg q$.
%   Persze átrendezve ekkor
%   \[
%   	\left( h\left( t \right)-h_1\left( t \right) \right)q\left( t \right)
%   	=
%   	r_1\left( t \right)-r\left( t \right)
%   \]
%   is fennáll.
%   Ekkor a fokszámokra figyelve
%   \[
%   	\deg\left( h-h_1 \right)+\deg q
%   	=
%   	\deg\left( r_1-r \right)
%   	\leq
%   	\max\left\{ \deg r_1,\deg (-r) \right\}
%   	<
%   	\deg q.
%   \]
%   Ez csak akkor lehetséges, ha $\deg\left( h-h_1 \right)=-\infty$,
%   ami azt jelenti, hogy $h=h_1$,
%   amiből persze $r_1=r$ már látszik is.
%end{proof}
}
\frame{
\begin{proposition}[a polinomgyűrű egy főideál-gyűrű]\label{pr:pgyurufoidealgyuru}
	A polinomok $\mathbb{F}\left[ t \right]$ kommutatív, egységelemes, nullosztómentes gyűrűjében
	minden a $\left\{ 0 \right\}$-tól különböző ideált generál az ideálban lévő egyetlen normált minimális fokszámú polinom.
	Így $\mathbb{F}\left[ t \right]$ egy főideál-gyűrű.
\end{proposition}
%begin{proof}
%   Legyen a $J$ egy ideálja $\mathbb{F}\left[ t \right]$-nek,
%   amely nem csak a zérus elemből áll.
%   Vegyünk egy minimális fokszámú de nem zérus polinomot $J$-ben,
%   tehát olyat,
%   amely maga sem zérus és nála kisebb fokszámú polinom már nincs $J$-ben a $0$ elemen kívül.
%   Legyen ez $d$.
%   Most megmutatjuk, hogy minden $p\in J$-re $d|p$.
%   A maradékos osztás szerint
%   valamely $h,r$ polinomokra
%   \[
%   	p\left( t \right)=
%   	h\left( t \right)d\left( t \right)+r\left( t \right);
%   	\text{ ahol }
%   	\deg r<\deg d.
%   \]
%   Mivel $p,d\in J$, és $J$ egy ideál, ezért $r\in J$.
%   No de, $d$ konstrukciója szerint ilyen csak a zérus polinom van,
%   ezért valóban $d|p$.
%   Ez éppen azt jelenti, hogy
%   \(
%   J=\left\{ dh:h\in\mathbb{F}\left[ t \right] \right\}
%   \)
%   azaz $d$ generálja a $J$ ideált.
%   Azt viszont már korábban is meggondoltuk,
%   hogy egy főideált csak egyetlen normált polinom generál.
%
%   Megmutattuk tehát,
%   hogy egyetlen normált, minimális fokszámú polinom van $J$-ben,
%   és minden $J$-beli polinom ennek többszöröse.
%end{proof}
Érdemes eltenni magunknak, hogy az ideál generáló eleme, tehát az ideálbeli elemek közös osztója
éppen az ideál minimális fokszámú nem zérus polinomja.
Ilyenből a normált polinomok közül csak egy van.
}
\end{document}
